\section{Introduction: one object, four registers}
\label{sec:intro}

This paper is the report of a single sustained experiment. We fix one mathematical
object --- a family of hyperbolic three-manifolds so small and so forced that it is
hard to call it a choice --- and we follow it wherever its own structure leads. It
leads, as it turns out, into four different parts of mathematics: the dynamics of a
trace map, the representation theory of a tower of Lie groups, the geometry of the
figure-eight knot's character variety, and the duality group of a four-dimensional
supersymmetric gauge theory. The thesis of the paper is not any one of the results
along the way. It is that these four faces are faces of \emph{the same object}, and
that following the object honestly --- recording at every step exactly what is proved,
what is computed, what is borrowed, and what is merely suggested --- maps out precisely
how far such an object reaches, and where it stops.

\subsection*{A note on proof status} Because the object touches so many registers, the
single greatest risk is a quiet drift in what is being claimed. We therefore attach to
every substantive statement a colour-coded badge:
\begin{center}\small
\sProved\ exact and machine-checkable\quad
\sExact\ symbolic over a ring\quad
\sNum\ high-precision numerics\\[2pt]
\sStruct\ structurally forced, not fully symbolic\quad
\sKnown\ established in the literature\\[2pt]
\sGated\ appears new but awaits a specialist\quad
\sPost\ proposed, not put through the verification gate
\end{center}
The badges are not decoration. The one genuinely new headline of the paper --- the
$\SL(4)$ figure-eight $A$-polynomial $L=-M^4$ and the family it suggests
(Section~\ref{sec:knot}) --- carries the badge \sGated\ throughout: an adversarial
literature read found no prior art, but an AI-assisted read \emph{de-risks} a novelty
question without \emph{closing} it, and the family is established only at two values of
the rank. We state this in the abstract, we state it again here, and we state it at the
result. Everything else in the paper is \sProved, \sExact, \sNum, or \sKnown, and is
labelled accordingly.

\subsection{The motivating question, and the firewall that contains it}

The object was not chosen for its knot theory. It was chosen as the smallest answer to
a deliberately na\"ive question: \emph{what is the least structure that is not nothing?}
One can refuse the metaphysics of ``why is there something'' and ask instead, as a
mathematical question, what the minimal non-degenerate self-referential object looks
like --- and discover that the minimal answers do not form a unique point but a
\emph{self-generating family}. That reframing is the project's philosophical spine, and
we will let it motivate the exposition. But motivation is all it is. We adopt one
absolute rule, which we will call the \emph{firewall}:

\begin{quote}\itshape
Philosophy and physical interpretation may cite mathematics; mathematics may never cite
philosophy. If a motivating sentence is ever needed to support a theorem, the boundary
has been crossed, and the sentence is wrong where it stands.
\end{quote}

This is why the paper can be, at once, philosophically motivated and mathematically
sober. The reader who wants only the theorems may read Sections~\ref{sec:object}--\ref{sec:chirality}
and ignore every interpretive remark without losing a single implication. The reader
interested in the larger question will find it addressed exactly once, in
Section~\ref{sec:bridge}, and addressed by the same standard of proof as everything
else: with a forced mathematical identification on one side and a confirmed obstruction
on the other.

\subsection{Architecture, not furniture}

A second principle organises what we do and do not claim. The object is a theory of
\emph{structure} --- of symmetry, of hierarchy, of which configurations are possible ---
and not a theory of \emph{values}: it does not select a coupling constant, a mass, or a
cosmological scale. Borrowing a phrase from the project's notes, it is architecture, not
furniture. This is not modesty for its own sake; it is a theorem-shaped fact, made
precise in Section~\ref{sec:bridge}, that the invariants the object naturally produces
are dimensionless and scale-free. Stating it at the outset pre-sorts the questions one
may ask: structural questions (what is the symmetry? what is the representation? what is
the $A$-polynomial?) have answers, and we give them; value questions (what number does
this predict?) meet a wall, and we exhibit the wall rather than climbing over it.

\subsection{What is in the paper}

Section~\ref{sec:object} introduces the metallic once-punctured-torus bundles and their
arithmetic. Section~\ref{sec:invariant} pins the boundary invariant $\kk=\tr[A,B]$ to
the Fricke--Vogt invariant of the Fibonacci trace map and identifies the metallic
monodromies with the $\SL(2,\Z)$ mapping-class action. Section~\ref{sec:tower} builds
the adjoint-type ``Dickson tower'' and proves its factorisation through $\SL(4)$, noting
the precise wall at $n\ge 5$. Section~\ref{sec:knot} is the geometric heart: the
trace-map fixed locus reproduces the figure-eight $A$-polynomial exactly at $\SL(2)$ and
$\SL(3)$, and yields the new, gated $\SL(4)$ relation $L=-M^4$, complete on the
irreducible locus. Section~\ref{sec:chirality} closes the chirality question with a
palindromic criterion for amphichirality. Section~\ref{sec:bridge} makes the bridge to
physics and confirms the wall. Sections~\ref{sec:discussion}--\ref{sec:open} discuss the
object's proper name as an aperiodic-order operator, the open problems, and the
reproducibility discipline; an appendix collects proof sketches.
